\subsection{Quantitative Results}
The set of human behavior research experiments presented in this work were designed to investigate three dimensions, Analytical vs Emotional Trust, Static vs. Interactive Conditions, and WEIRD vs. Normal participants. Across the six main quantitative hypotheses, and after Benjamini-Hochberg correction, our analysis found a robust pre- to post-experiment decrease in trust, no significant main effect of condition, and repeated differences between WEIRD and Normal participants in both trust level and trust updating pattern.

Our quantitative analysis began by establishing a significant effect of the LLM bias explanation exercise on the pre versus post experiment questionnaires. Across all participants, analytical, emotional, and overall trust each decreased significantly after the experiment, providing full support for Hypothesis 1. These results verify our experiment platform and design, as they demonstrate that the educational exercise produced a measurable change in AI trust responses. As expected, educating users on how AI can be biased reduced reported trust, while also showing that explanations of both AI bias and attempted bias mitigation can meaningfully reshape trust judgments.

After establishing this significant overall effect, our analysis continued on to comparing the effect of condition on changes in analytical, emotional, and overall trust. Contrary to our hypothesis, we did not observe a significant condition effect in relation to any of these trust outcomes, so Hypothesis 2 was not supported. This demonstrates that for our experiment platform, the interactivity of the educational material did not have a significant effect on changes in AI trust, and that the main effect of the intervention was not driven by whether the material was presented in interactive or static form.

After demonstrating no significant condition effect, our analysis turned to cultural group. For WEIRD participants alone, analytical trust and overall trust decreased significantly, but emotional trust did not remain significant after correction, so Hypothesis 3 received only partial support. For Normal participants, by contrast, analytical, emotional, and overall trust all decreased significantly, providing full support for Hypothesis 4. This difference suggests that while the intervention reduced trust in both cultural groups, the structure of that change was not identical, and emotional trust updating was less consistent in the WEIRD sample than in the Normal sample.

The final stage of our quantitative analysis evaluated direct contrasts between WEIRD and Normal participants. Hypothesis 5 received full support, with significant WEIRD-versus-Normal differences in analytical, emotional, and overall trust at both the pre- and post-experiment timepoints. Hypothesis 6 also found strong support for cultural differences in trust organization: the two groups differed significantly in both pre-trust and post-trust component contrasts, and Normal participants showed a significant analytical-versus-emotional change contrast, whereas the same within-group contrast was not significant for WEIRD participants. Taken together, these results support the central claim of this work that culture shapes not only how much people trust AI systems, but also how analytical and emotional trust are updated after learning about AI bias.